\section{Introduction}
\label{sec:intro}

High-dimensional approximate nearest-neighbour (ANN) search is a core primitive
in retrieval, recommendation, and vector databases. On a GPU, compressed search
is not governed by memory footprint alone. A representation can reduce bytes yet
increase lookup pressure, instruction issue, or irregular access; conversely, an
implementation can move the same logical bytes and still accelerate by issuing
fewer instructions. The systems problem is therefore to map the algorithm's
structure to the hardware bottleneck, not to assume that every reduction in
storage or every wider load is automatically beneficial.

JHQ~\cite{jhq} is a useful case study because its hierarchy exposes two
qualitatively different costs. After IVF routing, a compact primary code scores
many candidates, whereas a per-dimension residual code is read only for the
$c_k$ survivors of primary filtering. A naive port can mishandle both sides:
generic candidate-major codes give the warp a strided primary access pattern,
conventional 8-bit query tables grow with the number of subspaces, and an
over-large $c_k$ converts a selective residual stage into unnecessary work. JHQ
also leaves the refinement factor $\alpha$ in $c_k=\alpha k$ externally
specified. Thus the GPU problem has two coupled performance questions with two
different causes: how should the representation execute, and how much of the
expensive hierarchy should execute?

Our premise is that a GPU design should follow from explicit constraints and
rejected alternatives. We preserve the represented JHQ distances and ask, for
each hot-path decision, what resource is under pressure, what representation
property can relieve it, and what the alternative would cost. This separates
research claims from engineering controls: a transpose is useful because the
scan is candidate-parallel at fixed subspace, not because transposes are novel;
packing is useful because it removes issued work, not because it magically moves
fewer bytes; and a smaller LUT is useful only until extra lookups cost more than
additional residency saves.

The strongest representation-specific result comes from this last trade-off. A
conventional 8-bit ADC table materialises 256 distances per subspace. JHQ's
primary codebook is Cartesian, so squared Euclidean distance separates across
its one-dimensional levels. The same represented distance can therefore be
computed from two 16-entry tables. This is not simply ``smaller is better'': the
full table minimises lookups but competes for the same on-chip budget as the
scan's own candidate scratch, while a table-free formulation minimises table
storage but adds arithmetic and issued instructions. The two-table factorisation
is a middle point in this computation--residency design space. Mechanistically,
full-table state grows as $256M$ entries while the two-way factorised state
grows as $32M$, so the factorised representation approaches any fixed on-chip
capacity limit with an $8\times$ smaller coefficient. Experiments confirm the
prediction: the factorisation matters much more at $M{=}384$ than at $M{=}96$,
whereas still-smaller four- and eight-group tables are consistently slower.

The second contribution addresses a different decision: residual work. Refining
every routed candidate would erase the hierarchy's main advantage, while a fixed
small $\alpha$ can discard neighbours before the accurate stage sees them. We
therefore treat $c_k$ as a workload-level control variable. Our output-stability
criterion compares the returned top-$k$ set at a smaller budget with a generous
$\alpha_{\max}$ reference on sampled workload queries. We choose this signal
because it is available at deployment time without labelled nearest neighbours
or a learned query-difficulty model; the trade-off is that it is a proxy, not a
recall certificate, and its sample size must be chosen for the workload. A
matched CPU baseline shows this knob is not a portable one: pinning
$\alpha{=}100$ costs the CPU $8.7$--$41.8\%$ of its throughput and the GPU
$21.8$--$112.5\%$ on the same dataset and the same recall targets, so the
control JHQ leaves external is worth $2.5$ to $3\times$ more on the device we
port it to.

Our evaluation uses six datasets from 768 to 3{,}072 dimensions, up to 17.8M
vectors, on one RTX 5090. We use positive and negative controls to test the
design hypotheses rather than report only the winning implementation. Table-free
distance reduces table state yet loses $30$--$52\%$; four- and eight-way
factorisations are byte-identical to each other yet the eight-way form is
$5.8$--$38.3\%$ slower; packed loading moves the same logical bytes yet gains
$30$--$48\%$; a probe-cursor correction reduces issued loop work and can have
the largest speedup despite zero algorithmic novelty. These results jointly
indicate an issue-sensitive primary scan whose optimal design is not predicted
by byte counting alone. End-to-end, JHQ leads IVF-RaBitQ through
Recall@$10{=}0.95$ at large batch on four directly comparable datasets but
yields in parts of the high-recall tail, and batch size can change the relative
winner.

\smallskip\noindent\textbf{Contributions.}
First, we derive and evaluate a representation-driven GPU primary path for JHQ:
access-pattern-driven layout, representation-aligned packed loads, and an exact
Cartesian LUT factorisation whose two-way split is justified against full-table,
finer-factorisation, and table-free alternatives. The JHQ-specific research
claim is the factorisation and its hardware regime; standard layout and
instruction-level changes are supporting engineering.
Second, we propose a label-free, predictor-free workload policy for selecting
JHQ's residual refinement budget, quantify its sample risk and amortisation, and
show on a matched CPU baseline that the payoff of this control is several times
larger on the GPU than on the CPU.
Third, we evaluate the system with mechanism-oriented ablations, negative
results, recall--QPS frontiers, batch sensitivity, and build/memory costs,
explicitly reporting where competing GPU methods are stronger and where current
evidence does not generalise.
